
Dos barras se encuentran inicialmente a la misma temperatura t$_0$. Una de ellas tiene una longitud t$_{01}$=10.0 cm y un coeficiente de dilatación lineal $\alpha_1$ y la otra tiene una longitud t$_{02}$=12.0 cm, con un coeficiente de dilatación lineal $\alpha_2$. Se desea que al calentar las dos barras hasta una temperatura, la diferencia entre sus longitudes permanezca siempre igual a 2.0 cm, cualquiera que sea el valor de t. ¿Cuál debe ser el valor de la relación entre los coeficientes $\alpha_1$ y $\alpha_2$ para que esto suceda?